\documentclass[10pt,a4paper]{amsart}
\usepackage{amsmath,amssymb,amsthm,graphicx}
\usepackage{verbatim}
\newenvironment{reference}{\comment}{\endcomment}
%\newenvironment{reference}{}{}
\newenvironment{slogan}{\comment}{\endcomment}
\newenvironment{history}{\comment}{\endcomment}

% For commutative diagrams we use Xy-pic
\usepackage[all]{xy}

% We use 2cell for 2-commutative diagrams.
\xyoption{2cell}
% Theorem environments.
%
\theoremstyle{plain}
\newtheorem{theorem}[subsection]{Theorem}
\newtheorem{proposition}[subsection]{Proposition}
\newtheorem{lemma}[subsection]{Lemma}

\theoremstyle{definition}
\newtheorem{definition}[subsection]{Definition}
\newtheorem{example}[subsection]{Example}
\newtheorem{exercise}[subsection]{Exercise}
\newtheorem{situation}[subsection]{Situation}

\theoremstyle{remark}
\newtheorem{remark}[subsection]{Remark}
\newtheorem{remarks}[subsection]{Remarks}

\numberwithin{equation}{subsection}

% Macros
%
\def\lim{\mathop{\mathrm{lim}}\nolimits}
\def\colim{\mathop{\mathrm{colim}}\nolimits}
\def\Spec{\mathop{\mathrm{Spec}}}
\def\Hom{\mathop{\mathrm{Hom}}\nolimits}
\def\Ext{\mathop{\mathrm{Ext}}\nolimits}
\def\SheafHom{\mathop{\mathcal{H}\!\mathit{om}}\nolimits}
\def\SheafExt{\mathop{\mathcal{E}\!\mathit{xt}}\nolimits}
\def\Sch{\mathit{Sch}}
\def\Mor{\mathop{\mathrm{Mor}}\nolimits}
\def\Ob{\mathop{\mathrm{Ob}}\nolimits}
\def\Sh{\mathop{\mathit{Sh}}\nolimits}
\def\NL{\mathop{N\!L}\nolimits}
\def\CH{\mathop{\mathrm{CH}}\nolimits}
\def\proetale{{pro\text{-}\acute{e}tale}}
\def\etale{{\acute{e}tale}}
\def\QCoh{\mathit{QCoh}}
\def\Ker{\mathop{\mathrm{Ker}}}
\def\Im{\mathop{\mathrm{Im}}}
\def\Coker{\mathop{\mathrm{Coker}}}
\def\Coim{\mathop{\mathrm{Coim}}}

% Boxtimes
%
\DeclareMathSymbol{\boxtimes}{\mathbin}{AMSa}{"02}

%
% Macros for moduli stacks/spaces
%
\def\QCohstack{\mathcal{QC}\!\mathit{oh}}
\def\Cohstack{\mathcal{C}\!\mathit{oh}}
\def\Spacesstack{\mathcal{S}\!\mathit{paces}}
\def\Quotfunctor{\mathrm{Quot}}
\def\Hilbfunctor{\mathrm{Hilb}}
\def\Curvesstack{\mathcal{C}\!\mathit{urves}}
\def\Polarizedstack{\mathcal{P}\!\mathit{olarized}}
\def\Complexesstack{\mathcal{C}\!\mathit{omplexes}}
% \Pic is the operator that assigns to X its picard group, usage \Pic(X)
% \Picardstack_{X/B} denotes the Picard stack of X over B
% \Picardfunctor_{X/B} denotes the Picard functor of X over B
\def\Pic{\mathop{\mathrm{Pic}}\nolimits}
\def\Picardstack{\mathcal{P}\!\mathit{ic}}
\def\Picardfunctor{\mathrm{Pic}}
\def\Deformationcategory{\mathcal{D}\!\mathit{ef}}
\usepackage{cleveref}
\crefname{originalreference}{original reference}{original references}
\Crefname{originalreference}{Original reference}{Original references}
\makeatletter
\newlabel{stacks-morphisms-lemma-locally-finite-type-locally-noetherian}{{06R6}{source}{}{sourcepointer-1}{}}
\newlabel{stacks-morphisms-lemma-locally-finite-type-locally-noetherian@cref}{{[originalreference][0][]06R6}{[1][0][]source}{}{}{}}
\makeatother
\crefname{originalreference}{original reference}{original references}
\Crefname{originalreference}{Original reference}{Original references}
\makeatletter
\newlabel{stacks-morphisms-lemma-noetherian-finite-type-finite-presentation}{{0DQJ}{source}{}{sourcepointer-2}{}}
\newlabel{stacks-morphisms-lemma-noetherian-finite-type-finite-presentation@cref}{{[originalreference][0][]0DQJ}{[1][0][]source}{}{}{}}
\makeatother
\crefname{originalreference}{original reference}{original references}
\Crefname{originalreference}{Original reference}{Original references}
\makeatletter
\newlabel{stacks-sheaves-lemma-compare-quasi-coherent}{{0GQD}{source}{}{sourcepointer-3}{}}
\newlabel{stacks-sheaves-lemma-compare-quasi-coherent@cref}{{[originalreference][0][]0GQD}{[1][0][]source}{}{}{}}
\makeatother
\crefname{originalreference}{original reference}{original references}
\Crefname{originalreference}{Original reference}{Original references}
\makeatletter
\newlabel{sites-modules-lemma-local-pullback}{{03DO}{source}{}{sourcepointer-4}{}}
\newlabel{sites-modules-lemma-local-pullback@cref}{{[originalreference][0][]03DO}{[1][0][]source}{}{}{}}
\makeatother
\crefname{originalreference}{original reference}{original references}
\Crefname{originalreference}{Original reference}{Original references}
\makeatletter
\newlabel{spaces-cohomology-lemma-coherent-Noetherian}{{07UB}{source}{}{sourcepointer-5}{}}
\newlabel{spaces-cohomology-lemma-coherent-Noetherian@cref}{{[originalreference][0][]07UB}{[1][0][]source}{}{}{}}
\makeatother
\title{Four characterizations of coherent modules on locally Noetherian stacks}
\author{The Stacks Project Authors}
\date{stacks-project:stacks-cohomology.tex:local-snapshot-20261012; GNU FDL 1.2 or later}
\begin{document}
\maketitle
\noindent\textit{Editorial note.} All delivered mathematical text is the author's own, copied verbatim from the pinned Stacks Project chapter file stacks-cohomology.tex; no statement or proof step was rewritten, reordered or removed.\par\medskip
\noindent\textit{Editorial note.} The counted span cites the lemmas recorded in assembly.external\_references with canonical tags as navigation pointers.\par\medskip
\noindent\textit{Allowed standard prerequisites (editorial).}\begin{itemize}
\item Compare quasi-coherent restrictions. Sheaves on Stacks, Lemma stacks-sheaves-lemma-compare-quasi-coherent. x\textasciicircum{}*F as a pullback of the U-restriction.
\item Coherent modules on locally Noetherian algebraic spaces. Cohomology of Spaces, Lemma spaces-cohomology-lemma-coherent-Noetherian. Supports (1)=>(3).
\item Finite presentation vs finite type on ringed topoi. Modules on Sites lemmas; definition. (2)=>(1).
\end{itemize}\medskip
\ifdefined\hypertarget\hypertarget{sourcepointer-1}{}\fi
\noindent\textit{Source reference.} Original reference 06R6: sources/stacks-chow/tags.txt, source lines 8630--8630. Cross-chapter standard prerequisite cited by the author inside the counted span; it is used as a pointer and its own proof is not part of this item.\par\medskip
\ifdefined\hypertarget\hypertarget{sourcepointer-2}{}\fi
\noindent\textit{Source reference.} Original reference 0DQJ: sources/stacks-chow/tags.txt, source lines 17189--17189. Cross-chapter standard prerequisite cited by the author inside the counted span; it is used as a pointer and its own proof is not part of this item.\par\medskip
\ifdefined\hypertarget\hypertarget{sourcepointer-3}{}\fi
\noindent\textit{Source reference.} Original reference 0GQD: sources/stacks-chow/tags.txt, source lines 20858--20858. Cross-chapter standard prerequisite cited by the author inside the counted span; it is used as a pointer and its own proof is not part of this item.\par\medskip
\ifdefined\hypertarget\hypertarget{sourcepointer-4}{}\fi
\noindent\textit{Source reference.} Original reference 03DO: sources/stacks-chow/tags.txt, source lines 4427--4427. Cross-chapter standard prerequisite cited by the author inside the counted span; it is used as a pointer and its own proof is not part of this item.\par\medskip
\ifdefined\hypertarget\hypertarget{sourcepointer-5}{}\fi
\noindent\textit{Source reference.} Original reference 07UB: sources/stacks-chow/tags.txt, source lines 9966--9966. Cross-chapter standard prerequisite cited by the author inside the counted span; it is used as a pointer and its own proof is not part of this item.\par\medskip
\begin{lemma}
\label{lemma-coherent-Noetherian}
Let $\mathcal{X}$ be a locally Noetherian algebraic stack.
Let $\mathcal{F}$ be an $\mathcal{O}_\mathcal{X}$-module.
The following are equivalent
\begin{enumerate}
\item $\mathcal{F}$ is a quasi-coherent, finite type
$\mathcal{O}_\mathcal{X}$-module,
\item $\mathcal{F}$ is an $\mathcal{O}_\mathcal{X}$-module
of finite presentation,
\item $\mathcal{F}$ is quasi-coherent and for any morphism
$f : U \to \mathcal{X}$ where $U$ is a locally Noetherian algebraic space,
the pullback $f^*\mathcal{F}|_{U_\etale}$ is coherent, and
\item $\mathcal{F}$ is quasi-coherent and there exists an algebraic space
$U$ and a morphism $f : U \to \mathcal{X}$ which is locally of finite type,
flat, and surjective, such that the pullback $f^*\mathcal{F}|_{U_\etale}$
is coherent.
\end{enumerate}
\end{lemma}

\begin{proof}
Let $f : U \to \mathcal{X}$ be as in (4).
Then $U$ is locally Noetherian (Morphisms of Stacks, Lemma
\ref{stacks-morphisms-lemma-locally-finite-type-locally-noetherian})
and we see that the statement of the lemma makes sense. Additionally,
$f$ is locally of finite presentation by Morphisms of Stacks, Lemma
\ref{stacks-morphisms-lemma-noetherian-finite-type-finite-presentation}.
Let $x$ be an object of $\mathcal{X}$ lying over the scheme $V$.
In order to prove (2) we have to show that, after replacing $V$
by the members of an fppf covering of $V$, the restriction $x^*\mathcal{F}$
has a global finite presentation on $\mathcal{X}/x \cong (\Sch/V)_{fppf}$.
The projection $W = U \times_\mathcal{X} V \to V$ is
locally of finite presentation, flat, and surjective.
Hence we may replace $V$ by the members of an \'etale covering of $W$
by schemes and assume we have a morphism $h : V \to U$ with $f \circ h = x$.
Since $\mathcal{F}$ is quasi-coherent, we see that
the restriction $x^*\mathcal{F}$
is the pullback of $h_{small}^*(f^*\mathcal{F})|_{U_\etale}$
by $\pi_V$, see Sheaves on Stacks, Lemma
\ref{stacks-sheaves-lemma-compare-quasi-coherent}.
Since $f^*\mathcal{F}|_{U_\etale}$ locally in the
\'etale topology has a finite presentation by assumption, we conclude
(4) $\Rightarrow$ (2).

\medskip\noindent
Part (2) implies (1) for any ringed topos (immediate from the definition).
The properties ``finite type'' and ``quasi-coherent''
are preserved under pullback by any morphism of ringed topoi, see
Modules on Sites, Lemma \ref{sites-modules-lemma-local-pullback}.
Hence (1) implies (3), see Cohomology of Spaces, Lemma
\ref{spaces-cohomology-lemma-coherent-Noetherian}.
Finally, (3) trivially implies (4).
\end{proof}

\end{document}
